% Figure: laminar pipe velocity profiles for power-law fluids. The Newtonian
% parabola (n = 1) sits between the blunt shear-thinning profile (n < 1, more
% plug-like) and the pointed shear-thickening profile (n > 1). Plotted as the
% normalised velocity u/u_max against radius r/R with the pipe axis at r = 0.
% Profile: u/u_max = 1 - (r/R)^{(n+1)/n}.
\begin{figure}[htbp]
\centering
\begin{tikzpicture}
\begin{axis}[
  width=11cm, height=7cm,
  xlabel={radius $r/R$ (wall at $1$, axis at $0$)},
  ylabel={velocity $u/u_{\max}$},
  xmin=0, xmax=1, ymin=0, ymax=1.08,
  legend pos=south west,
  legend cell align=left,
  grid=both, grid style={enggray!18},
  minor grid style={enggray!8},
  tick label style={font=\scriptsize},
  label style={font=\small},
  legend style={font=\scriptsize},
]
% Shear-thinning n = 0.6, exponent (n+1)/n = 2.667 -> blunt.
\addplot[engblue, very thick, domain=0:1, samples=80] {1 - x^2.667};
\addlegendentry{shear-thinning $n=0.6$ (blunt)}
% Newtonian n = 1, exponent 2 -> parabola.
\addplot[engred, very thick, domain=0:1, samples=80] {1 - x^2};
\addlegendentry{Newtonian $n=1$ (parabola)}
% Shear-thickening n = 1.5, exponent (n+1)/n = 1.667 -> pointed.
\addplot[engamber, very thick, domain=0:1, samples=80] {1 - x^1.667};
\addlegendentry{shear-thickening $n=1.5$ (pointed)}
\end{axis}
\end{tikzpicture}
\caption[Power-law pipe velocity profiles]{Laminar pipe profiles for a
power-law fluid $\tau = K(du/dr)^n$. The velocity is $u/u_{\max} = 1 -
(r/R)^{(n+1)/n}$, so the flow index $n$ sets the shape. The Newtonian parabola
($n=1$) is the middle curve; a shear-thinning fluid ($n<1$) carries a blunter,
more plug-like profile that puts proportionally more of its flow near the axis,
raising the mean-to-centreline ratio above the Newtonian one-half; a
shear-thickening fluid ($n>1$) carries a more pointed profile. Worked
example~22 integrates the $n=0.6$ curve to the power-law flow-rate law.}
\label{fig:vol03ch12:power-law-profiles}
\end{figure}
